\chapter{The destination}

This curriculum treats Latin as a language to be read and written, not as a
code in which each word is mechanically replaced by an English word. Its
principal destination is the \Missal: you should eventually be able to open
the book at an unfamiliar Mass, identify the grammatical spine of each text,
follow its prayer, argument, narrative, or poem directly in Latin, and consult
a lexicon only for genuinely unfamiliar vocabulary or names.

That destination requires wider Latin, not narrower memorization. A person who
can read only familiar formulas has not yet gained transferable reading
ability. The course therefore begins with transparent ordinary sentences,
moves through varied prose and poetry, and returns continually to the Missal
with stronger grammar, vocabulary, and literary judgment.

\section{What broad competence includes}

The common trunk develops the ability to read:

\begin{itemize}
  \item project-composed Latin about ordinary objects, actions, places, and
        human relations;
  \item biblical narrative, exhortation, prayer, and poetry;
  \item selected classical prose, including increasingly demanding Ciceronian
        examples and periods;
  \item patristic and ecclesiastical prose, especially Augustine and other
        excellent models appropriate to the available grammar;
  \item selected theological, canonical, rubrical, and historical registers;
  \item classical, biblical, Christian, hymnic, and liturgical poetry at the
        level needed by the common course;
  \item the Ordinary, Canon, prefaces, Proper of Time, Proper of Saints,
        Commons, votive and Requiem texts, and principal Holy Week material in
        the identified Missal; and
  \item complete unfamiliar formularies without a facing translation.
\end{itemize}

Composition belongs to the same formation. It begins with changing one form,
continues through phrases, transparent sentences, clause combination, and
controlled imitation, and culminates in source-aware advanced prose. Writing
reveals whether forms and constructions are available from memory rather than
merely recognizable on the page.

\section{What the course does not add}

Pronunciation, phonetics, singing, chant notation, and ceremonial instruction
are outside the present course. Scansion in the poetry branch concerns textual
quantity and verse structure; it does not become a pronunciation guide.
Grammar does not decide liturgical use, doctrine, textual authority, or canon
law by itself.

\chapter{How to study without a calendar}

There are no weeks, daily quotas, uniform sessions, or completion dates. The
packet order preserves dependencies, but the amount of time and repetition
belongs to the material and the person doing the work. Move quickly through
what is genuinely transparent; remain with what is not yet dependable.

\section{Definition and retrieval}

A genuinely new technical term is defined where it is first taught. The
Reference Grammar gathers those decisions in one retrieval home. Later
packets use established terms without restating them unless a new distinction
or a narrower meaning matters.

\Definition{working definition}{A concise statement of the features needed to
recognize and use a concept at the present stage. A later lesson may refine it
but may not silently contradict it.}

\Definition{qualitative readiness}{A description of what you should reasonably
be able to recognize, explain, read, or compose before continuing. It is
tested on fresh Latin and names the exact work worth repeating when
uncertainty recurs.}

\Definition{retrieval round}{A fresh attempt to recall vocabulary, forms,
rules, or constructions without copying the answer. Short and delayed rounds
make the material increasingly available in connected reading and writing.}

When a definition still feels abstract, use a four-part notebook entry:

\begin{enumerate}
  \item copy the definition exactly;
  \item underline the feature that distinguishes it from its nearest rival;
  \item supply one Latin example and label the decisive evidence; and
  \item write one non-example or contrast and explain why it differs.
\end{enumerate}

Knowing a label is not enough. A term has become useful when you can identify
it in fresh Latin, cite the evidence, distinguish a neighboring construction,
and use it accurately where composition permits.

\section{The study loop}

\StudyLoop

Handwriting is encouraged because it makes endings, word order, and revision
visible. Typed work can supplement it. What matters is that the first decision
remain available for comparison with the correction.

\section{The error record}

Keep a notebook or loose-leaf record divided into five families:

\begin{enumerate}
  \item forms and dictionary entries;
  \item agreement and government;
  \item clause boundaries and syntax;
  \item vocabulary, idiom, and register; and
  \item composition, rhetoric, and revision.
\end{enumerate}

For a substantive correction, record the original decision, the corrected
form or analysis, the evidence that decides between them, and a new example.
Recurring uncertainty points to an explanation or retrieval family. It does
not call for restarting the course or repeating unrelated pages.

\chapter{The common trunk}

The twenty-two trunk packets group forty-six established lessons. Packet IDs
give the new substantial publication sequence; the stable lesson IDs remain
visible inside it and continue to identify teaching, exercises, passages, and
answers.

\section{Foundations: M01--M06}

Foundations begins with ordinary project-composed Latin: a fox is brown, a
girl carries water, friends see a road. Concrete subject matter keeps the first
forms transparent. Missal phrases and short source excerpts enter as soon as
the grammar can carry them without turning recognition of a familiar prayer
into a substitute for reading.

\begin{xltabular}{\linewidth}{@{}P{0.55in}P{1.0in}P{2.0in}Y@{}}
  \toprule
  \textbf{Packet} & \textbf{Stable lessons} & \textbf{Center} & \textbf{Before continuing} \\
  \midrule\endhead
  M01 & EL-I.1--I.2 & Form, function, case, number, gender, agreement, and the
  English-grammar bridge & Map a fresh simple clause, distinguish form from
  job, and justify basic case and agreement from Latin evidence \\
  M02 & EL-I.3--I.4 & First, second, and third declensions & Recover stems,
  decline representative nouns, resolve common ambiguous endings in context,
  and use them in ordinary sentences \\
  M03 & EL-I.5--I.6 & Fourth and fifth declensions; adjective systems and
  agreement & Control the full noun system and connect adjectives with their
  heads across different-looking endings \\
  M04 & EL-I.7--I.8 & \Latin{sum}, \Latin{possum}, and the present active
  system & Recall the forms readily, distinguish predicates from objects, and
  narrate or describe simple present, past, and future actions \\
  M05 & EL-I.9--I.10 & Perfect active, passive systems, and deponents & Track
  time, voice, agency, participial agreement, and deponent form and sense in a
  fresh paragraph \\
  M06 & EL-I.11--I.12 & Pronouns, prepositions, reference, and a repeatable
  reading method & Follow reference and government through connected Latin,
  read a short unfamiliar passage by method, and write a coherent paragraph
  from controlled vocabulary \\
  \bottomrule
\end{xltabular}

Composition moves from substitution and inflection in M01 to connected
ordinary prose in M06. The vocabulary strand begins with concrete objects,
colors, people, places, and actions before adding frequent biblical,
ecclesiastical, and Missal vocabulary.

\section{Core grammar: M07--M12}

Core grammar develops the non-finite and clause structures needed for mature
prose. Readings broaden deliberately: narrative and reported speech, selected
classical sentences, Vulgate prose, and early patristic passages stand beside
liturgical examples.

\begin{xltabular}{\linewidth}{@{}P{0.55in}P{1.1in}P{2.0in}Y@{}}
  \toprule
  \textbf{Packet} & \textbf{Stable lessons} & \textbf{Center} & \textbf{Before continuing} \\
  \midrule\endhead
  M07 & EL-II.1--II.2 & Relative, interrogative, and indefinite reference;
  comparison, adverbs, and number & Follow reference chains and degree through
  a fresh passage and use the same resources in clear description \\
  M08 & EL-II.3--II.4 & Infinitives, indirect statement, participles, and the
  ablative absolute & Identify controllers, subjects, voice, relative time,
  and attachment before translating or recomposing \\
  M09 & EL-II.5--II.7 & Subjunctive morphology and sequence; purpose, result,
  indirect command, and substantive clauses & Form the subjunctive system,
  classify clauses from converging evidence, and compose connected examples
  without treating \Latin{ut} as a complete explanation \\
  M10 & EL-II.8--II.10 & Temporal, causal, concessive, conditional, and
  characteristic relations & Explain speaker stance and clause relation in a
  mixed paragraph, then vary those relations in composition \\
  M11 & EL-II.11--II.12 & Commands, prohibitions, wishes, fearing, gerund,
  gerundive, supine, and periphrastics & Distinguish volition, fear,
  obligation, purpose, and verbal-noun functions in reading and writing \\
  M12 & EL-II.13--II.14 & Place, time, means, manner, cause, marked order,
  ellipsis, attraction, and compression & Reconstruct a difficult sentence
  without losing printed forms and explain what its actual order contributes \\
  \bottomrule
\end{xltabular}

By M12, controlled writing has become multi-clause composition. Ciceronian and
Augustinian examples can now be studied for architecture as well as decoded
for propositions, while every attribution and adaptation remains explicit.

\section{Missal-centered reading: M13--M18}

These packets reorganize the accumulated grammar by genre. They remain broad
Latin packets: biblical, patristic, classical, and ecclesiastical parallels
test transfer and illuminate genuine differences of register. The Missal now
becomes the sustained center rather than the only source of sentences.

\begin{xltabular}{\linewidth}{@{}P{0.55in}P{1.1in}P{2.0in}Y@{}}
  \toprule
  \textbf{Packet} & \textbf{Stable lessons} & \textbf{Center} & \textbf{Before continuing} \\
  \midrule\endhead
  M13 & EL-III.1--III.2 & Independent lookup discipline and the Ordinary &
  Read recurring formulas as grammar rather than remembered translation and
  carry the same method into an unrelated prose passage \\
  M14 & EL-III.3--III.4 & Collects, Secrets, Postcommunions, compressed
  petition, gift, and effect & Map unfamiliar orations without forcing one
  label where Latin permits more than one relation; imitate a supplied frame \\
  M15 & EL-III.5--III.6 & Lessons, Epistles, Gospels, argument, narrative,
  speech, and biblical idiom & Follow a complete argument or narrative,
  distinguish its register, summarize directly, and recast part in simpler
  Latin \\
  M16 & EL-III.7--III.8 & Chants, sequences, prefaces, and the Roman Canon &
  Recover poetic and periodic dependencies, mark responsible ellipsis, and
  explain what prose order would lose \\
  M17 & EL-III.9--III.10 & Temporal, Sanctoral, and Commons & Transfer
  vocabulary and genre knowledge across seasons, saints, and reusable
  formulary patterns without guessing from familiarity \\
  M18 & EL-III.11--III.12 & Holy Week, Requiem, and complete unfamiliar Masses
  & Sustain direct reading across every genre of a fresh formulary, log only
  genuine lookups, and explain consequential grammatical decisions \\
  \bottomrule
\end{xltabular}

Poetry in M16 and M18 is treated as poetry: word displacement, parallelism,
ellipsis, imagery, verbal recurrence, and written-ending pattern are part of
reading. The neutral prose scaffold is a tool, not a claim that poetic order
is ornamental or defective.

\section{Common advanced trunk: M19--M22}

The advanced trunk brings excellent prose models to the foreground. Cicero,
Augustine, biblical and patristic prose, liturgical periods, theology, canon
law, and other audited selections teach register, rhetoric, source awareness,
imitation, and sustained composition.

\begin{xltabular}{\linewidth}{@{}P{0.55in}P{1.1in}P{2.0in}Y@{}}
  \toprule
  \textbf{Packet} & \textbf{Stable lessons} & \textbf{Center} & \textbf{Before continuing} \\
  \midrule\endhead
  M19 & EL-A1--A2 & Register across time, Vulgate texture, and biblical idiom &
  Compare exact passages, describe real differences without caricature, and
  write short imitations in two declared registers \\
  M20 & EL-A3--A4 & Allusion, context, rhetoric, and periodic prose & Trace or
  reject a proposed source responsibly, diagram a long period, and explain how
  exact rhetorical members order thought \\
  M21 & EL-A5--A6 & Textual collation, commentary, and controlled imitation &
  Keep witness, normalization, analysis, and editorial choice distinct while
  changing propositions inside an identified stylistic frame \\
  M22 & EL-A7--A8 & Guided and independent composition, revision, and research &
  Plan, draft, revise, annotate, and defend substantial Latin prose whose
  grammar, idiom, register, and source use remain accountable \\
  \bottomrule
\end{xltabular}

Advanced composition is therefore the culmination of a strand begun in M01,
not an activity attached after reading has ended. M22 also supplies the common
advanced abilities assumed by the poetry branch.

\chapter{The advanced poetry branch}

The branch begins after the common advanced trunk and consists of eight
separately printable packets. It deepens poetry without displacing advanced
prose or implying that the common trunk omitted poetic Missal texts.

\begin{xltabular}{\linewidth}{@{}P{0.55in}P{2.0in}Y@{}}
  \toprule
  \textbf{Packet} & \textbf{Center} & \textbf{Developed ability} \\
  \midrule\endhead
  P1 & Textual prosody & Mark quantity, feet, elision, caesura, and poetic
  syntax on the page while keeping textual analysis distinct from
  pronunciation \\
  P2 & Hexameter & Scan and interpret selected elevated narrative, especially
  Virgilian models, relating syntax, line movement, and imagery \\
  P3 & Elegiac verse & Read the couplet's contrast and closure, compressed
  diction, and selected classical transformations \\
  P4 & Lyric meters & Recognize selected Horatian architectures and explain
  how stanza, syntax, and rhetorical progression cooperate \\
  P5 & Christian poetics & Compare classical inheritance, biblical language,
  and Christian transformation in the verified Pentecost Sequence and
  clearly labeled project comparisons \\
  P6 & Rhythmic hymnody and sequences & Distinguish quantity from syllable
  count, line grouping, repeated syntax, and written endings while reading the
  complete supplied sequence and project hymn-style comparisons as literary
  wholes, without inferring pronunciation \\
  P7 & Poetic commentary & Write source-aware commentary on diction, allusion,
  imagery, sound, syntax, meter, and theological or liturgical context without
  confusing grammatical analysis with disciplinary authority \\
  P8 & Verse composition & Move from substitution and line imitation through
  constrained stanzas to revised original verse with a declared model and a
  source-and-idiom record \\
  \bottomrule
\end{xltabular}

The opening of each poetry packet states the forms, syntax, and literary
analysis it assumes. A person may return to a trunk reading, rhetoric, or
composition packet whenever those foundations need renewal.

\chapter{From translation to direct reading}

Translation remains useful when it exposes grammatical relations. It becomes
less central as direct comprehension grows.

\begin{enumerate}
  \item \textbf{Foundations and core grammar:} write structural English that
        keeps Latin relations visible, then natural English.
  \item \textbf{Early sustained reading:} annotate clause and dependency
        structure and give a concise English summary rather than replacing
        every line.
  \item \textbf{Later Missal reading:} paraphrase transparent clauses in
        easier Latin and answer comprehension questions from the Latin.
  \item \textbf{Advanced work:} use English chiefly for grammatical, textual,
        and literary commentary; understand and discuss the text from its
        Latin form.
\end{enumerate}

Never hide uncertainty beneath elegant English. During study, an awkward
structural version that reveals the syntax is more useful than polished prose
built on a mistaken subject or clause relation.

\chapter{Vocabulary, memory, and reference discipline}

The vocabulary sequence overlaps six broad strands: ordinary concrete Latin;
general high-frequency Latin; classical prose; biblical and patristic Latin;
theological and canonical Latin; and liturgical and Missal-specific formulas.
A word may belong to several strands. The controlling registry records its
citation form, working senses, morphology, usage or government, strand, first
appearance, and later returns.

Each packet prints compact new, short-return, longer-return, and cumulative
memory rounds. The Memory Workbook supplies fuller detachable sheets for
Latin-to-English recognition, English-to-Latin production, dictionary forms,
principal parts, paradigms, usage and government, grammar facts,
reconstruction, and controlled composition. Use blank forms for
additional rounds without replacing the populated course sequence.

Look up a word only after recording a morphological hypothesis: lemma, part of
speech, likely form, and likely function. Then choose a sense that fits the
construction, author, and context. Copying the first gloss is not reading.

For an unfamiliar passage, mark words as known directly, inferable from form
or family, or genuinely needing lookup. The final ability permits lookup for
rare names and uncommon vocabulary. It does not rely on a facing translation
or a parser that makes grammatical decisions in advance.

\chapter{Cumulative self-review}

Four separate companions gather fresh work after broad stretches of the
trunk: Foundations, Core Grammar, Missal Reading, and Advanced Reading and
Composition. Parallel forms permit a later return without simply remembering
the first answer. The Missal-reading companion prints two identified
formularies for its first returns and a fixed fresh pair for transfer; no
outside Missal or source-selection procedure is needed.

Use these reviews diagnostically. Compare answers by their deciding evidence,
record recurring uncertainty, revisit the named teaching or memory family,
and return later with different Latin. A sound result is described in
abilities: forms recalled in fresh contexts, clauses explained from evidence,
connected Latin understood directly, and composition revised by defensible
rules and models.

Fluency is retained access. While beginning a new packet, continue a small
amount of earlier vocabulary retrieval, copying, parsing, direct reading, and
composition so that previous forms remain available in new contexts.
